\section{What the Boundary Costs}
\label{sec:ch12-what-the-boundary-costs}
\index{boundary!what stops being checkable|(}%

Chapter~\ref{ch:entity-resolution-done-right} ended with a rule and handed it
here: a field that crosses a subgraph boundary should be nullable unless you can
guarantee referential integrity across a network, and you almost certainly
cannot. It said this chapter would apply that five more times.

It applies three times, and the three it does not are the more useful half.

\subsection*{Three references, all nullable}

\index{nullability!a reference to somebody else's entity}%
These three are the arrows from section~\ref{sec:ch12-the-order-the-graph-decides}
and they all became nullable:

\begin{minted}[fontsize=\footnotesize]{text}
Review.author       Customer!  ->  Customer
Order.customer      Customer!  ->  Customer
OrderLine.product   Product!   ->  Product
\end{minted}

\code{OrderLine.product} is the one
chapter~\ref{ch:the-first-cut-extracting-catalog} got wrong in a comment and
chapter~\ref{ch:entity-resolution-done-right} measured. That comment claimed a
dangling identifier
\enquote{surfaces as a null product at the router}; it does not, because the
field was \code{Product!} inside \code{lines: [OrderLine!]!}, so a product
Catalog cannot resolve takes the whole order history with it. The only reason it
never bit is that the seeder builds every line from a product it has just
written, and a guarantee that holds because of how the test data was made is not
a guarantee.

The other two are new boundaries and got the same treatment for the same reason.
What is worth looking at is what the resolvers became, because the shape is the
argument:

\begin{minted}[fontsize=\footnotesize,breaklines]{csharp}
public static Customer? GetAuthor([Parent] Review review)
    => new() { Id = review.CustomerId };
\end{minted}

Fourteen lines became two. The old version called a DataLoader, got a whole
customer out of the same database, and threw if there was not one. This one
wraps a key it is already holding.

\index{measurement!work that moved rather than went away}%
Two things follow and they pull in opposite directions. It got cheaper here: no
DataLoader, no batch, no statement, so a page of a hundred and twenty reviews
costs Reviews nothing at all for its authors, which is the statement that left
in section~\ref{sec:ch12-the-plan-grows-a-parallel-node}. And it got more
expensive somewhere else: the router makes a fifth fetch, and the twelve
distinct customers behind those reviews arrive at Accounts as twelve
representations. The work did not go away. It moved, and grew a network hop on
the way.

It also lost something that is easy to miss. Inside the monolith,
\code{Review.author} and \code{Order.customer} called the same DataLoader
instance, so a query that reached a customer through their orders and again
through their reviews fetched them once and handed both fields the same object.
Two processes, two request scopes, two caches. What replaces it is the router's
own batching, and whether those two sets of keys end up in the same entity fetch
is the planner's decision rather than mine.
Chapter~\ref{ch:inside-the-router-and-the-composer} reads the planner properly
and owns the answer.

\subsection*{Where the rule stops, and why}

\index{nullability!contributed fields are not references}%
Three of the six boundary fields did not change, and I want to make the argument
rather than let it look like an oversight, because
chapter~\ref{ch:entity-resolution-done-right} predicted five.

\begin{minted}[fontsize=\footnotesize]{text}
Product.price              Money!  unchanged
Product.shippingCost       Money!  unchanged
Product.availableQuantity  Int!    unchanged
\end{minted}

The difference is between a reference and a contribution. \code{Review.author}
is a reference: this service hands the router a key and the router has to go and
find the thing, and finding may fail. \code{Product.price} is not a reference at
all. The router has already located the entity; Pricing is being asked a
question about a key it has been handed, and either it has the row or it does
not, and its answer is authoritative either way.

Which means the two failures are different failures. A dangling reference is
ordinary in a distributed system, and a schema that pretends otherwise is
lying. A product Catalog has and Pricing does not is a broken invariant between
two databases, and answering null for it would hand a client a hole where a
price should be, with no way to tell that from a product that is genuinely
free. Pricing throws instead, which
chapter~\ref{ch:entity-resolution-done-right} measured turning into
\code{Unexpected Execution Error} at the boundary. Loud and unhelpful is the
right trade for a state that should not exist.

\code{Product.availableQuantity} is the third and it is neither: Inventory
answers zero for a product it has never counted, because \enquote{we have never
counted this} is not something a storefront can render, and the field is total.

\index{nullability!the rule with its boundary}%
So the rule chapter~\ref{ch:entity-resolution-done-right} gave has a boundary of
its own, and it is worth stating
in the sharper form: \textbf{a reference to an entity you do not own must be
nullable. A field you contribute to an entity somebody else owns need not be,
and the test is whether the failure is a missing datum or a broken invariant.}
That is more useful than making everything nullable, which is where a rule of
thumb about boundaries eventually lands if nobody argues with it.

\subsection*{The check that could not be given back}

\index{submitReview@\code{submitReview}!loses its second check}%
Chapter~\ref{ch:the-first-cut-extracting-catalog} took
\code{ProductNotFoundError} off \code{submitReview}, because Reviews could no
longer ask whether a product existed, and recorded that this chapter would argue
about whether \code{@requires} could give it back. Here is the argument, and
also the same thing happening a second time.

Accounts is a separate service now, so Reviews cannot check a customer either.
\code{@requires} cannot help. It feeds a field resolver on an entity the router
has already located, and the mutation is a root field handed a customer
identifier by a client that may have invented it. There is no entity to locate
and no subgraph the router would think to consult.

That leaves three options and none of them is good. A synchronous call to
Accounts on the write path gives the check back and puts a second service's
availability in front of every review anybody submits. Leaving
\code{CustomerNotFoundError} in the union with nothing able to raise it is the
worst of the three, because it is a branch a client will write and never reach.
So the third: accept the write.

\begin{graphqlsdl}[fontsize=\footnotesize]
union SubmitReviewError = RatingOutOfRangeError | DuplicateReviewError
\end{graphqlsdl}

Four members at chapter~\ref{ch:schema-design-that-survives-change}, three at
chapter~\ref{ch:the-first-cut-extracting-catalog}, two now. The collection asserts
what that means rather than describing it. A review submitted against a
\code{Customer} node identifier nobody has ever seeded succeeds, the payload's
errors are null, and the code that used to name the problem is nowhere in the
response, because nothing can raise it.

\index{errors!a rule you own against data you do not}%
What survives is the pair of rules Reviews owns outright. A rating is between
one and five, and one customer reviews one product once, and both still raise a
typed error. The line between what left and what stayed is exact: a domain that
owns a rule can still enforce it, and a domain that has stopped owning the data
cannot.

Chapter~\ref{ch:the-first-cut-extracting-catalog} called its version of this a
decision. Doing it a second time is
closer to a consequence, and that is the honest summary of what a split costs at
the boundary. Every check you had that reached across a seam is now a check
somebody else could do and you cannot, and the schema is where you admit it.
\code{Review.author} answering null at read time is the same check, moved from
write to read and from an error a client can handle to a hole it has to.

\medskip
Which is the whole balance sheet, for the storefront query the gate measures:
four statements before and four after, two extra resolvers, two extra round
trips, three fields that admit they might not answer, one error type gone, and a
monolith that is not in the repository any more. Ask for the reviews as well and
it is three extra round trips rather than two, which is the shape of the whole
trade: the bill grows with how far across the graph a query reaches, not with
how many services there are. Chapter~\ref{ch:hard-modeling-problems} takes the arrangement this
chapter produced and asks the questions it cannot answer: what a page of
something looks like when the pages are on two services, what \code{Query.node}
means now, and what to do about the enum that is declared twice.

\index{boundary!what stops being checkable|)}%
